%%
% 摘要信息
% 本文档中前缀"c-"代表中文版字段, 前缀"e-"代表英文版字段
% 摘要内容应概括地反映出本论文的主要内容，主要说明本论文的研究目的、内容、方法、成果和结论。要突出本论文的创造性成果或新见解，不要与引言相 混淆。语言力求精练、准确，以 300—500 字为宜。
% 在摘要的下方另起一行，注明本文的关键词（3—5 个）。关键词是供检索用的主题词条，应采用能覆盖论文主要内容的通用技术词条(参照相应的技术术语 标准)。按词条的外延层次排列（外延大的排在前面）。摘要与关键词应在同一页。
% modifier: 黄俊杰(huangjj27, 349373001dc@gmail.com)
% update date: 2017-04-15
%%

% 中文
\cabstract{
根据实验内容 \cite{lab01}, 在此次实验需要完成个人所得税计算的程序应用，此报告参照评分标准\cite{labeval01}，完成这份实验报告撰写并附上相应的实验报告说明。
本实验实现了一个基于命令行的个人所得税计算器，采用 Java 语言并遵循面向对象设计方法。系统将税率档位与起征点抽象为可配置模块，支持动态调整、输入校验与异常处理，并配备 JUnit 单元测试以覆盖正常、边界与异常情形。通过对示例（如月收入 4300 元、起征点 1600 元）进行分段超额累进税率计算，验证了计算结果的正确性与程序的健壮性。本实验既强化了封装、抽象与异常处理等编程思想，也为将来扩展多税制支持或图形界面提供了良好基础。
}
\ckeywords{个人所得税；Java；面向对象；异常处理；单元测试}


% \ckeywords{} % 保持仅英文摘要时，避免模板读取中文关键词元数据报错
% % 英文
% \eabstract{
% After learning the lecture one of compiler principles \cite{lecture01}, we are supposed to finish the homework.
% }
% \ekeywords{Compiler Principles, JAVA}

